%% ============================================================
%% Academic Journal Paper Template
%% Compatible with: IEEE, ACM, Elsevier, Springer conventions
%% ============================================================

\documentclass[10pt]{article}

%% ---- Core Packages ----------------------------------------
\usepackage[T1]{fontenc}
\usepackage[utf8]{inputenc}
\usepackage{lmodern}
\usepackage{microtype}           % Microtypographic refinements

%% ---- Page Geometry ----------------------------------------
\usepackage[
  top=1in, bottom=1in,
  left=0.75in, right=0.75in,
  columnsep=0.25in
]{geometry}

%% ---- Mathematics ------------------------------------------
\usepackage{amsmath, amssymb, amsthm}
\usepackage{mathtools}
\usepackage{bm}                  % Bold math symbols
\usepackage{filecontents}

%% ---- Figures & Tables -------------------------------------
\usepackage{graphicx}
\usepackage{booktabs}            % Professional table rules
\usepackage{multirow}
\usepackage{array}
\usepackage{caption}
\usepackage{subcaption}
\usepackage{float}

%% ---- Code Listings ----------------------------------------
\usepackage{listings}
\usepackage{xcolor}
\usepackage{amsmath}
\usepackage{algorithm}
\usepackage{algpseudocode}

\lstdefinestyle{codestyle}{
  backgroundcolor=\color{gray!8},
  basicstyle=\ttfamily\footnotesize,
  breakatwhitespace=false,
  breaklines=true,
  captionpos=b,
  commentstyle=\color{green!50!black},
  keywordstyle=\color{blue}\bfseries,
  stringstyle=\color{orange!70!black},
  numberstyle=\tiny\color{gray},
  numbers=left,
  numbersep=5pt,
  showstringspaces=false,
  frame=single,
  rulecolor=\color{gray!40},
  tabsize=2
}
\lstset{style=codestyle}

%% ---- Hyperlinks -------------------------------------------
\usepackage[
  colorlinks=true,
  linkcolor=blue!70!black,
  citecolor=green!50!black,
  urlcolor=blue!60!black
]{hyperref}
\usepackage{url}
\usepackage{arrayjobx}
%% ---- Bibliography -----------------------------------------
\usepackage[numbers, sort&compress]{natbib}

%% ---- Miscellaneous ----------------------------------------
\usepackage{lipsum}              % Lorem ipsum filler — remove in production
\usepackage{enumitem}
\usepackage{siunitx}             % SI units: \SI{3.14}{\mega\hertz}
\usepackage{cleveref}            % Smart cross-references: \cref{fig:foo}

%% ---- Theorem Environments ---------------------------------
\theoremstyle{plain}
\newtheorem{theorem}{Theorem}[section]
\newtheorem{lemma}[theorem]{Lemma}
\newtheorem{corollary}[theorem]{Corollary}
\newtheorem{proposition}[theorem]{Proposition}

\theoremstyle{definition}
\newtheorem{definition}[theorem]{Definition}
\newtheorem{example}[theorem]{Example}

\theoremstyle{remark}
\newtheorem{remark}[theorem]{Remark}

%% ---- Custom Commands --------------------------------------
\newcommand{\R}{\mathbb{R}}
\newcommand{\N}{\mathbb{N}}
\newcommand{\Z}{\mathbb{Z}}
\newcommand{\E}{\mathbb{E}}
\newcommand{\Prob}{\mathbb{P}}
\newcommand{\norm}[1]{\left\lVert #1 \right\rVert}
\newcommand{\abs}[1]{\left\lvert #1 \right\rvert}
\newcommand{\inner}[2]{\left\langle #1,\, #2 \right\rangle}
\newcommand{\ie}{\textit{i.e.}\xspace}
\newcommand{\eg}{\textit{e.g.}\xspace}
\newcommand{\etal}{\textit{et al.}\xspace}

%% ---- Caption Formatting -----------------------------------
\captionsetup{
  font=small,
  labelfont=bf,
  format=hang,
  justification=justified
}

%% ---- Header / Footer --------------------------------------
\usepackage{fancyhdr}
\pagestyle{fancy}
\fancyhf{}
\fancyhead[R]{\small\thepage}
\renewcommand{\headrulewidth}{0.4pt}

%% ============================================================
%%  FRONT MATTER
%% ============================================================

\title{%
  \vspace{-1.5em}%
  \large\textbf{Part Four}%
}

\author{%
  \textbf{A. E. Mahrouss}$^{1}$
}

\date{%
  \small \today
}

%% ============================================================
\begin{document}
%% ============================================================

%% ============================================================

\begin{definition}
	Let:
	\begin{align*}
		E(x) &= \frac{\gamma(x)}{\delta(x) + \lambda(x)}; \quad \gamma(x) \neq 0.
	\end{align*}
	And. Let based on $E(x)$:
	\begin{align*}
		P' &= \prod_{n=1}^{3}\rho(x) x^{n} + \int E(x) dx;\\
		P' &= \prod_{n=1}^{3}\rho(x) x^{n} + \int \frac{\gamma(x)}{\delta(x) + \lambda(x)} dx; \quad P(s) \neq 0;\\
		P' &= \int \frac{\gamma(x)}{\delta(x) + \lambda(x)} dx + x^{n} \rho(x)^{3}.
	\end{align*}
	And. The following for $x$:
	\begin{align*}
		\prod_{n=1}^{3}\rho(x) x^{n} + \int \frac{\gamma(x)}{\delta(x) + \lambda(x)} dx; \quad P(s) \neq 0.
	\end{align*}
	Admits a solution in $\mathbb{R}$ of the following:
	\begin{align*}
		P'' &= \prod_{n=1}^{3}\rho(x) x^{n} + \int \frac{\gamma(x)}{\delta(x) + \lambda(x)} dx; \quad P'' \neq 0.
	\end{align*}
	Separately, let the following:
	\begin{align*}
		L' + \lambda(x)^{2} &= \frac{\gamma(x)}{\delta(x)};\\
	 	L' + \lambda(x) &= \frac{\gamma(x)}{\delta(x)  \lambda(x) };\\
	 	L' &= \frac{\gamma(x)}{\delta(x)  \lambda(x) } - \lambda(x); \quad \delta(x)\lambda(x) \neq 0.
	\end{align*}
	Where we must find a value of $\ln x$ that equals:
	\begin{align*}
		&= \frac{\ln x}{\Delta y};\\
		\Delta y &= \ln x;
	\end{align*}
	Separately, we must find $I'$ that equals:
	\begin{align*}
		r &= ( \frac{\ln x}{\Delta n} + \frac{\Delta n}{\ln x})^m; \quad r \neq 0; \quad r \in \mathbb{R}{+};\\
		I' &= - r + ( \frac{\ln x}{\Delta n} + \frac{\Delta n}{\ln x})^m; \quad r \neq 0; \quad r \in \mathbb{R}{+};
	\end{align*}
	And. For all variables of $r$ being equal to each other, e.g: $x =n = \ln x = m.$ Then the following:
	\begin{align*}
		r &= ( \frac{\ln x}{n} + \frac{n}{\ln x})^m; \quad r \neq 0; \quad r \in \mathbb{R}{+}.
	\end{align*}
	Admits a transecdenctal number.
\end{definition}

\begin{theorem}
	The following plots both a real and imiginary part for both:
	\begin{align*}
		\frac{7}{\log n} + \frac{363 \log 7}{140} &= \sum_{n=1}^{m} \frac{\ln x}{n} + \frac{n}{\ln x}; \quad m = 7
	\end{align*}
	And:
	\begin{align*}
		\frac{7}{\log n} + \frac{\log^{7} 7}{5040} &= \prod_{n=1}^{m} \frac{\ln x}{n} + \frac{n}{\ln x}; \quad m = 7.
	\end{align*}
\end{theorem}

\nocite{*}

\bibliographystyle{acm}


\end{document}
